\subsection*{Two retained failed choices at the remaining norm-seventeen model}

\textbf{Exploratory certificate and geometric interpretation; not promoted.}
The uniform degree-seventeen theorem remains open.  On
\(C_0:y^2+y=x^5\), with its Weierstrass point \(W\), two selected
polarized self-isogenies of multiplier seventeen have nonnodal theta images.
This is an obstruction to these choices, not to all multiplier-seventeen
isogenies or to the desired monodromy theorem.

The first choice is \(3+2g+2h\), with anticommuting order-four
automorphisms.  The second is
\(1+2g_1+2g_2+2g_3+2g_4\), with four pairwise anticommuting
order-four automorphisms.  The latter units are explicitly retained in
\texttt{post\_CP30\_seventeen\_four\_unit\_interpolation.json}.
The cross terms in the Rosati product cancel, giving norm
\(9+4+4=17\) and \(1+4+4+4+4=17\), respectively.

For both choices, exact maximal minors prove that the prescribed divisor
of degree seven or nine imposes independent conditions in
\(H^0(9W)\) or \(H^0(11W)\), at every finite parameter \(p=x(P)\).
The resulting residual divisor \(E_P\) has degree two and is unique;
therefore it is nonspecial.  Its class is the negative of the isogeny
image of \(P-W\).  In particular, no finite point maps to zero.
The retained norm computations give its Mumford coordinates
\[
 u_P(x)=x^2+L(p)x+M(p),\qquad
 v_P(x)=\nu(p)x+\kappa(P).
\]
All finite coordinate denominators are supported on the displayed pole
polynomial.  For the second map, that polynomial has degree twelve but
eleven distinct roots.  The reduced pole algebra has a residual-x
generator of degree eleven with separable minimal polynomial; hence its
pole parameters have distinct residual-x values.  The repeated factor
was retained and handled on its reduced support, not discarded as a
simple-root assumption.

The polynomial-combination certificates for the nonpole coordinate
equations, followed by exact quotient-field greatest common divisors,
produce the following packets over the field with binary modulus
\(\texttt{0x10eb}\).  All arrays in this paragraph encode coefficients
in that fixed field.  For the first map, let \(p\) be a root of
\([2623,1821,1]\).  The common off-diagonal polynomial in \(q\) is
\[
\begin{split}
 q^4+372q^3+(3362+372p)q^2
 +(373+373p)q+(2622+2622p).
\end{split}
\]
For the second map, let \(p\) be a root of
\([1821,372,1]\).  The corresponding polynomial is
\[
 q^2+(2890+2890p)q+(3158+2623p).
\]
The exact checker verifies that these polynomials are squarefree, have
no root \(q=p\), avoid the entire pole support, and divide all three
coordinate-equality polynomials for \(L,M,\nu\).  It also verifies
\(L\ne0\).  Thus the two maps have respectively five and three
distinct finite x parameters with common \((L,M,\nu)\).

Here is the geometric passage from these coordinates to actual fibres.
For fixed \(L,M,\nu\), reduce
\((\nu x+\kappa)^2+(\nu x+\kappa)-x^5\) modulo
\(x^2+Lx+M\).  The coefficient of \(x\) is independent of
\(\kappa\), while the constant coefficient is an equation
\(\kappa^2+\kappa=c\).  Since one residual divisor exists, the
coefficient-of-x condition holds; the constant equation has exactly two
solutions over the algebraically closed field.  Replacing \(P\) by
its hyperelliptic conjugate changes every residual ordinate by one,
and hence exchanges these two solutions.  At each of the distinct
x parameters one can therefore choose the unique ordinate giving
the same \(\kappa\).  The resulting actual residual divisors, and
therefore the isogeny images, agree.  Because \(L\ne0\), the residual
quadratic is separable.  These are actual multi-point fibres, rather
than solely a large resultant bad set.

All original interpolation, rational-coordinate, polynomial-combination
and quotient-field outputs remain unchanged.  The final reduced-root
check is
\texttt{post\_CP30\_seventeen\_C0\_packet\_obstructions.py},
with its same-stem JSON output.  It passed on
20 November 2026 at 01:15:14.999896 UTC in 0.0203053080040263 seconds.
This note has not received an independent full proof audit and is
excluded from the numbered proved paper.
